\section{RARE vs Limiting RARE}
\label{sec:lim-rare-equiv}

Recall the definition of \limRARE{} from \cref{sec:lim-rare}.

\thmRareEquivDefn*
\begin{proof}
Fix the allocation rule $A$.
(We omit its superscript in this proof.)
First, suppose that $\pvec$ is $\delta$-additive $\kappa$-RARE.
For every seller $j \in [m]$ and deviation $q_j \ge 0$, there exists a response
$\qvec_{-j} \in \IR_j(q_j, \pvec_{-j})$ such that
\[ D_j(q_j, \pvec_{-j}) \le R_j(\qvec) \le \kappa R_j(\pvec) + \delta. \]
Taking the supremum over $q_j \ge 0$ gives
$V_j(\pvec_{-j}) \le \kappa R_j(\pvec) + \delta$ for every seller $j$.
Thus, $\pvec$ is $\delta$-additive $\kappa$-\limRARE{}, proving the last claim.

Now suppose that $\pvec$ is $\delta$-additive $\kappa$-\limRARE{}.
Fix any $\eps > 0$, seller $j \in [m]$, and deviation $q_j \ge 0$.
The set $\IR_j(q_j, \pvec_{-j})$ is nonempty since it contains $\pvec_{-j}$,
and revenues are nonnegative, so its defining infimum $D_j(q_j, \pvec_{-j})$ is finite.
By the definition of infimum, there exists a response
$\qvec_{-j} \in \IR_j(q_j, \pvec_{-j})$ such that
\[ R_j(\qvec) < D_j(q_j, \pvec_{-j}) + \eps
    \le V_j(\pvec_{-j}) + \eps
    \le \kappa R_j(\pvec) + \delta + \eps. \]
Hence, $\pvec$ is $(\delta + \eps)$-additive $\kappa$-RARE.

Conversely, suppose that $\pvec$ is $(\delta + \eps)$-additive $\kappa$-RARE
for every $\eps > 0$.
By the first part of the proof, for every seller $j \in [m]$ and every $\eps > 0$,
\[ V_j(\pvec_{-j}) \le \kappa R_j(\pvec) + \delta + \eps. \]
Letting $\eps$ tend to zero gives $V_j(\pvec_{-j}) \le \kappa R_j(\pvec) + \delta$,
so $\pvec$ is $\delta$-additive $\kappa$-\limRARE{}.
\end{proof}
